\documentclass[UTF8]{ctexart}
\usepackage{xeCJK}
\usepackage{geometry}
\geometry{left=2cm,right=2cm,top=2.5cm,bottom=2.5cm}
\setlength{\headheight}{12.7pt}

\usepackage{titlesec}
\titleformat{\section}
  {\normalfont\large\bfseries}
  {\thesection}
  {1em}
  {}
\titlespacing*{\section}
  {0pt}
  {3.5ex plus 1ex minus .2ex}
  {2.3ex plus .2ex}

\usepackage{amsmath}
\usepackage{amssymb}
\usepackage{amsthm}
\usepackage{enumitem}
\setlist[enumerate]{itemsep=0pt,topsep=0pt,parsep=0pt,leftmargin=0pt,itemindent=2.5em}
\setlist[itemize]{itemsep=0pt,topsep=0pt,parsep=0pt,leftmargin=0pt,itemindent=2.5em}
\usepackage{fancyhdr}
\pagestyle{fancy}
\fancyhf{}
\usepackage{graphicx}
\usepackage{hyperref}
\usepackage{indentfirst}
\usepackage{lmodern}

\begin{document}
\setlength{\abovedisplayskip}{1pt}
\setlength{\belowdisplayskip}{1pt}

\section{一阶理论与模型}

\hypertarget{sec:定义1.1}{}
\noindent\textbf{定义1.1}

给定一阶语言$\mathcal{L}$, $\mathcal{L}$-结构$A$和语句$\varphi$.

如果存在$A$上的变元赋值$\sigma$使得$(A,\sigma)\models\varphi$,
就称结构$A$满足语句$\varphi$, 记作$A\models\varphi$.

\vspace{10pt}

设$T$是一族语句, 用$A\models T$表示$A$满足$T$中的所有语句.

\vspace{10pt}

\hypertarget{sec:定理1.1}{}
\noindent\textbf{定理1.1}

如果$A\models\varphi$, 那么对$A$上任意的变元赋值$\tau$都有$(A,\tau)\models\varphi$.

\noindent{\kaishu 证明.}

依定义存在$\sigma$使得$(A,\sigma)\models\varphi$. 因为$\varphi$没有自由变元,
所以根据\href{01 重合引理#nameddest=sec:定理1.2}{重合引理}有
$(A,\tau)\models\varphi$. \hfill $\qedsymbol$

\vspace{10pt}

设$T$是一族语句, $\varphi$是语句, 则显然$T\models\varphi$当且仅当对任意的结构$A$
有$A\models T\to A\models\varphi$.

\vspace{10pt}

\hypertarget{sec:定义1.2}{}
\noindent\textbf{定义1.2 (一阶理论)}

给定一阶语言$\mathcal{L}$, 称其中的一族语句$T$为一个$\mathcal{L}$-\textbf{理论}.

\vspace{10pt}

\hypertarget{sec:定义1.3}{}
\noindent\textbf{定义1.3 (一阶模型)}

给定一阶语言$\mathcal{L}$和$\mathcal{L}$-理论$T$, 称满足$T$的结构为$T$的\textbf{模型}.

\vspace{10pt}

任给一阶语言$\mathcal{L}$和$\mathcal{L}$-结构$A$,
把$\mathcal{L}$中所有被$A$满足的语句的全体记为
\[
    \mathrm{Th}_{\mathcal{L}}(A)=\{\varphi\in F^{\mathcal{L}}\mid
    \varphi\text{\ 是语句\ },A\text{\ 满足\ }\varphi\},
\]
称之为$A$的理论. 显然$A$自身是该理论的模型.





\newpage

\section{理论的完备性}

\hypertarget{sec:定义2.1}{}
\noindent\textbf{定义2.1 (理论的完备性)}

给定一阶语言$\mathcal{L}$和$\mathcal{L}$-理论$T$. 如果对任意的语句$\varphi$有
\[
    T\models\varphi\quad\lor\quad T\models\neg\varphi,
\]
就称$T$是\textbf{完备的}.

\vspace{10pt}

\hypertarget{sec:例2.1}{}
\noindent\textbf{例2.1}

任给一阶语言$\mathcal{L}$和$\mathcal{L}$-结构$A$,
$A$的理论$\mathrm{Th}_{\mathcal{L}}(A)$一定是完备的.

\noindent{\kaishu 证明.}

对任意的语句$\varphi$:
\begin{enumerate}
    \item 如果$A\models\varphi$, 则$\varphi\in\mathrm{Th}_{\mathcal{L}}(A)$,
    从而$\mathrm{Th}_{\mathcal{L}}(A)\models\varphi$,
    \item 如果$A\models\neg\varphi$,
    同理$\mathrm{Th}_{\mathcal{L}}(A)\models\neg\varphi$. \hfill $\qedsymbol$
\end{enumerate}

\vspace{10pt}

\hypertarget{sec:定理2.1}{}
\noindent\textbf{定理2.1}

给定一阶语言$\mathcal{L}$. 对$\mathcal{L}$的任意的一致理论$T$,
$T$是完备的当且仅当$T$的所有模型都有相同的理论, 或者说$T$的所有模
型\href{../04 Löwenheim-Skolem 定理/02 初等等价#nameddest=sec:定义1.1}{初等等价}.

\noindent{\kaishu 证明.}

\noindent{\kaishu 必要性.}

假设$T$完备, 任取$T$的模型$A,B$. 对任意的语句$\varphi$, 不妨设$A\models\varphi$.

此时$T\cup\{\varphi\}$一致, 所以$T\not\models\neg\varphi$,
由完备性可知$T\models\varphi$, 从而$B\models\varphi$.

\noindent{\kaishu 充分性.}

假设$T$的所有模型初等等价, 任取一个模型$A$. 对任意的语句$\varphi$:
\begin{enumerate}
    \item 如果$A\models\varphi$, 那么$T$的模型都满足$\varphi$, 即$T\models\varphi$,
    \item 如果$A\models\neg\varphi$, 同理$T\models\neg\varphi$. \hfill $\qedsymbol$
\end{enumerate}

\vspace{10pt}

\hypertarget{sec:定理2.2}{}
\noindent\textbf{定理2.2}

给定一阶语言$\mathcal{L}$和$\mathcal{L}$-结构$A,B$.
如果$B\models\mathrm{Th}_{\mathcal{L}}(A)$, 那么$A,B$初等等价.

\noindent{\kaishu 证明.}

因为$\mathrm{Th}_{\mathcal{L}}(A)$完备而$A,B$都是它的模型. \hfill $\qedsymbol$

\end{document}
